% !TEX program = xelatex

\documentclass[12pt,a4paper]{article}

\usepackage{fontspec}
\usepackage{polyglossia}
\setmainlanguage{russian}
\setotherlanguage{english}
\setmainfont{DejaVu Serif}
\newfontfamily\cyrillicfonttt{DejaVu Sans Mono}
\usepackage{amsmath,amssymb,amsfonts,amsthm}
\usepackage{geometry}
\usepackage{booktabs}
\usepackage{tabularx}
\usepackage{titlesec}
\usepackage{hyperref}
\usepackage{cite}
\usepackage{microtype}
\microtypesetup{expansion=false}
\setlength{\emergencystretch}{3em}

\titleformat{\section}{\large\bfseries}{\thesection}{1em}{}
\titleformat{\subsection}{\normalsize\bfseries}{\thesubsection}{1em}{}

\geometry{
    left=25mm,
    right=20mm,
    top=25mm,
    bottom=25mm
}

\hypersetup{
    colorlinks=true,
    linkcolor=blue,
    citecolor=blue,
    urlcolor=blue
}

\begin{document}

\section{Введение}

Адсорбция полиэлектролитов~\cite{Hara1993, Schmitz1993} представляет значительный практический интерес~\cite{Kawaguchi1992, Fleer1993, Bajpai1997}. Данное явление находит применение в различных областях науки и техники, таких как коллоидная химия, медицина, фармация, пищевая промышленность, очистка воды и др. За последние 30~лет было приложено немало экспериментальных и теоретических усилий для выяснения факторов, определяющих адсорбцию полиэлектролитов на заряженных поверхностях~\cite{Kawaguchi1992, Fleer1993, Bajpai1997}.

Большинство теоретических работ, посвященных адсорбции полиэлектролитов на заряженной поверхности, выполнено в рамках метода самосогласованного поля (ССП)~\cite{Fleer1993}. В этих теориях распределение плотности полимера связано с локальным электростатическим потенциалом через комбинацию уравнения Пуассона--Больцмана и уравнения Эдвардса, описывающего конформации полимера в эффективном внешнем потенциале (см. обзоры~\cite{Fleer1993, Bajpai1997}). Данный подход был впервые применен Ван дер Схее и Ликлемой~\cite{VanderSchee1984}, а также Эверсом и соавт.~\cite{Evers1986}. Они показали, что в бессолевых растворах сильное отталкивание между заряженными мономерами приводит к формированию очень тонких адсорбционных слоев. При экранировании этого взаимодействия добавлением соли количество адсорбированного полимера возрастает, а адсорбционный слой становится толще~\cite{Blaakmeer1990}. Обобщение теории Ван дер Схее и Ликлемы на случай слабых полиэлектролитов было выполнено Бёмером и соавт.~\cite{Bohmer1990}.

Адсорбция полиэлектролитов также исследовалась с использованием приближения доминирования основного состояния метода ССП~\cite{Varoqui1991, Varoqui1993, Borukhov1995, Chatellier1996}. Линеаризованные решения уравнений Пуассона--Больцмана и диффузии были получены Вароки и соавт.~\cite{Varoqui1991}. Они рассмотрели конформацию слабозаряженных полиэлектролитов на границе раздела жидкость--твердое тело и рассчитали изотерму адсорбции, а также профиль концентрации полимеров вблизи заряженной границы. Численное решение нелинейного уравнения Пуассона--Больцмана было представлено Боруховым и соавт.~\cite{Borukhov1995}. Эти авторы рассчитали профиль концентрации слабозаряженных полиэлектролитов между двумя заряженными поверхностями.

Несмотря на значительные успехи метода ССП в выявлении ключевых факторов, управляющих адсорбцией полиэлектролитов~\cite{Fleer1993}, данный метод опирается на численное решение системы самосогласованных уравнений. Это крайне затрудняет разделение эффектов различных факторов и не позволяет дать простую физическую интерпретацию полученных результатов. В настоящей работе будет показано, что возможно получить аналитическое решение уравнений самосогласованного поля в случае, когда равновесная концентрация полиэлектролита вблизи заряженной поверхности определяется балансом между электростатическим притяжением заряженных мономеров к поверхности и взаимодействиями исключенного объема. В частности, для случая $\theta$-растворителя этот баланс приводит к параболическому профилю плотности полиэлектролита в адсорбированном слое. 

Статья организована следующим образом. Раздел~II посвящен решению нелинейного уравнения Пуассона--Больцмана для адсорбции полиэлектролита в $\theta$-растворителе. В разделах~III и~IV полученные результаты обобщаются на случай плохого и хорошего растворителей.

\section{Структура адсорбированного слоя вблизи заряженной поверхности в $\theta$-растворителе}

\subsection{Модель}

Рассмотрим адсорбцию полиэлектролитных цепей со степенью полимеризации $N$ и долей $f$ отрицательных зарядов, случайно распределенных вдоль цепи, на положительно заряженной поверхности с поверхностной плотностью заряда $e\sigma$. Поверхность расположена параллельно плоскости $xy$ в точке $z = 0$. Цепи адсорбируются из разбавленного раствора с диэлектрической проницаемостью $\varepsilon$; средняя концентрация полимера и концентрация соли задаются величинами $\rho_{\text{bulk}}$ и $\rho_{\text{salt}}$ соответственно. Диссоциированные противионы полимерных цепей распределены в растворе со средней плотностью $f\rho_{\text{bulk}}$. Вдали от заряженной поверхности конформации полимерных цепей не возмущены электростатическими взаимодействиями с противоположно заряженной поверхностью, а конформация полиэлектролитной цепи контролируется кулоновским отталкиванием между заряженными мономерами. Это электростатическое отталкивание приводит к вытягиванию цепи на масштабах длины, превышающих размер электростатического блоба~\cite{deGennes1979, Dobrynin1995} $D_e$. 

Конформация цепи внутри электростатического блоба практически не возмущена электростатическими взаимодействиями, при этом число мономеров в нем составляет $g_e \approx (D_e / a)^2$ для случая $\theta$-растворителя для остова полимера. Размер электростатического блоба $D_e$, содержащего $g_e$ мономеров, может быть найден из условия равенства электростатической энергии блоба $e^2 g_e^2 f^2 / (\varepsilon D_e)$ тепловой энергии $kT$. Это дает размер электростатического блоба $D_e \approx a u^{-1/3} f^{-2/3}$ и число мономеров в нем $g_e \approx u^{-2/3} f^{-4/3}$ (где $u = l_B / a$ --- отношение длины Бьёррума $l_B = e^2 / (\varepsilon k T)$ к размеру связи $a$). В разбавленном растворе конфигурация полиэлектролита представляет собой вытянутую цепь из $N/g_e$ электростатических блобов длиной
\begin{equation}
L \approx \frac{N}{g_e} D_e \approx a u^{1/3} f^{2/3} N.
\end{equation}

Электростатическая собственная энергия полиэлектролитной цепи $W_{\text{pe}}$ пропорциональна тепловой энергии $kT$, умноженной на число электростатических блобов $N/g_e$ в цепи:
\begin{equation}
W_{\text{pe}} \approx k T u^{2/3} f^{4/3} N \approx k T \frac{N}{g_e}.
\end{equation}
Объединяя эту энергию с вкладом трансляционных степеней свободы цепей, можно получить химический потенциал полиэлектролитной цепи вдали от заряженной поверхности:
\begin{equation}
\mu_{\text{ch}} \approx k T \ln\left( \frac{\rho_{\text{bulk}} a^3}{N} \right) + k T \frac{N}{g_e}.
\end{equation}

Вблизи заряженной поверхности полиэлектролитные цепи образуют концентрированный раствор полимера. В этой области распределение плотности полимера может быть описано в рамках приближения среднего поля в предположении, что электростатические взаимодействия полиэлектролита с эффективным полем, создаваемым другими цепями, преобладают над его собственной электростатической энергией. В этом приближении плотность полимера $\rho(z)$ и плотность малых ионов $\rho_\alpha(z)$ зависят только от расстояния $z$ до заряженной поверхности. Данное приближение справедливо до тех пор, пока локальная плотность полимера $\rho(z)$ превышает плотность внутри электростатического блоба $\rho_e \approx g_e / D_e^3 \approx a^{-3} u^{1/3} f^{2/3}$.

В рамках приближения среднего поля свободная энергия системы содержит три слагаемых~\cite{Varoqui1991, Varoqui1993, Borukhov1995, Chatellier1996, Borukhov1998, Joanny1999}: электростатические взаимодействия между противионами, полиэлектролитными цепями и поверхностными зарядами $V_{\text{elect}}$, трансляционную энтропию противионов $F_{\text{count}}$ и свободную энергию полимерных цепей $F_{\text{pol}}$:
\begin{equation}
\Delta F = V_{\text{elect}} + F_{\text{pol}} + F_{\text{count}}.
\end{equation}
Часть свободной энергии, описывающая электростатические взаимодействия поверхностных зарядов с флуктуациями плотности раствора, $V_{\text{elect}}$, имеет вид:
\begin{equation}
\frac{V_{\text{elect}}}{k T} = -2\pi l_B \sigma S \int_0^\infty z q(z) dz - \pi l_B S \int_0^\infty \int_0^\infty |z - z'| q(z) q(z') dz dz'.
\end{equation}
Первое слагаемое в правой части описывает эффекты электростатического взаимодействия между заряженной поверхностью и распределением заряда $q(z)$ в системе:
\begin{equation}
q(z) = -f \rho(z) + \sum_\alpha q_\alpha \rho_\alpha(z),
\end{equation}
где $q_\alpha$ и $\rho_\alpha(z)$ --- валентность и локальная плотность малых ионов в системе, $\rho(z)$ --- локальная концентрация полимера. Второе слагаемое в правой части формулы~(5) описывает электростатические взаимодействия между заряженными слоями, расположенными на расстояниях $z$ и $z'$ от заряженной поверхности. Полимерный вклад $F_{\text{pol}}$ в свободную энергию включает три слагаемых: конформационную энтропию полимерных цепей, третий вириальный член, описывающий взаимодействия мономер--мономер в $\theta$-растворителе, и идеально-газовый вклад трансляционной энтропии цепей в адсорбированном слое:
\begin{equation}
\frac{1}{S}\frac{F_{\text{pol}}}{k T} = \frac{a^2}{6} \int_0^D \left( \frac{d\sqrt{\rho(z)}}{dz} \right)^2 dz + \frac{a^6}{3} \int_0^D \rho(z)^3 dz + \frac{\Gamma}{N} \ln\left( \frac{\Gamma a^3}{D N e} \right),
\end{equation}
где $\Gamma$ --- степень заполнения поверхности полимером:
\begin{equation}
\Gamma = \int_0^D \rho(z) dz.
\end{equation}
Полиэлектролитные цепи локализованы внутри слоя толщиной $D$ и находятся в тепловом равновесии с объемом раствора. В состоянии равновесия химический потенциал полиэлектролитных цепей в адсорбированном слое
\begin{equation}
\mu_{\text{ads}} = k T \ln\left( \frac{\Gamma a^3}{D N} \right) + N\mu,
\end{equation}
где $\mu$ --- множитель Лагранжа для ограничения~(8), равен химическому потенциалу в объеме.

Неэлектростатический вклад малых ионов в свободную энергию~(4) учитывается на уровне их идеально-газовой энтропии:
\begin{equation}
\frac{1}{S}\frac{F_{\text{count}}}{k T} = \sum_\alpha \int_0^\infty \rho_\alpha(z) \ln\left( \frac{\rho_\alpha(z)}{e} \right) dz.
\end{equation}

Для нахождения равновесных распределений плотности полимера и противионов вблизи заряженной поверхности необходимо взять вариационную производную свободной энергии по плотностям полимера $\rho(z)$ и малых ионов $\rho_\alpha(z)$ с учетом дополнительных связей, фиксирующих полное число молекул в системе. Экстремальное уравнение для распределения плотности полимера принимает вид:
\begin{equation}
\frac{\mu}{k T} = a^6 \rho(z)^2 - f \psi(z) - \frac{a^2}{6}\frac{1}{\sqrt{\rho(z)}}\frac{d^2\sqrt{\rho(z)}}{dz^2}.
\end{equation}
Параметр $\mu$ может быть найден из условия равенства химических потенциалов цепей в объеме $\mu_{\text{ch}}$ (уравнение~(3)) и в адсорбированном слое $\mu_{\text{ads}}$ (уравнение~(9)). Для длинных цепей ($N \gg 1$) параметр $\mu$ равен $k T g_e^{-1}$ и может рассматриваться как химический потенциал мономера в объеме.

В уравнении~(11) введен приведенный электростатический потенциал $\psi(z)$:
\begin{equation}
\psi(z) = -2\pi l_B \sigma z - 2\pi l_B \int_0^\infty |z - z'| q(z') dz',
\end{equation}
удовлетворяющий уравнению Пуассона
\begin{equation}
\frac{d^2\psi(z)}{dz^2} = 4\pi l_B \left( f \rho(z) - \sum_\alpha q_\alpha \rho_\alpha(z) \right)
\end{equation}
с граничным условием на заряженной поверхности:
\begin{equation}
\left.\frac{d\psi(z)}{dz}\right|_{z=0} = -4\pi l_B \sigma.
\end{equation}
Предполагается, что специфические взаимодействия между полимерами и поверхностью отсутствуют (кроме электростатических), а сама поверхность непроницаема для мономеров. Это дает следующее граничное условие для концентрации полимера на поверхности:
\begin{equation}
\rho(0) = 0.
\end{equation}
Распределение малых ионов в адсорбированном слое подчиняется больцмановскому распределению:
\begin{equation}
\rho_\alpha(z) = \rho_\alpha \exp(-q_\alpha \psi(z)),
\end{equation}
где $\rho_\alpha$ --- объемная концентрация малых ионов. Таким образом, система уравнений~(11)--(13) совместно с граничными условиями~(14) и~(15) полностью описывает адсорбцию полиэлектролита из разбавленного раствора.

\subsection{Автомодельный адсорбированный слой}

Система дифференциальных уравнений~(11)--(13) может быть решена аналитически, когда экранирование заряженной поверхности определяется адсорбированными полиэлектролитными цепями и ионами соли. В случае слабого электростатического потенциала $\psi(z) \ll 1$ эта система сводится к:
\begin{equation}
\frac{\mu}{k T} = a^6 \rho(z)^2 - f \psi(z) - \frac{a^2}{6}\frac{1}{\sqrt{\rho(z)}}\frac{d^2\sqrt{\rho(z)}}{dz^2},
\end{equation}
\begin{equation}
\frac{d^2\psi(z)}{dz^2} - \frac{\psi(z)}{r_D^2} = 4\pi l_B f \rho(z),
\end{equation}
где $r_D$ --- дебаевский радиус экранирования ионами соли ($r_D^{-2} = 8\pi l_B \rho_{\text{salt}}$). При анализе данных уравнений сначала пренебрежем последним слагаемым в правой части уравнения~(17), связанным с конформационной энтропией цепей, а также граничным условием на заряженной поверхности~(15). Границы применимости этого приближения будут рассмотрены ниже. Для длинных полиэлектролитных цепей ($N \gg 1$) множитель Лагранжа $\mu \sim k T g_e^{-1}$, и им можно пренебречь, если локальная концентрация мономеров $\rho(z)$ больше $a^{-3} u^{1/3} f^{2/3}$. С учетом этого получаем выражение для плотности полимера~\cite{deGennes1976, deGennes1979, Dobrynin1995, Israelachvili1985}:
\begin{equation}
a^6 \rho(z)^2 \approx f \psi(z),
\end{equation}
что позволяет переписать уравнение~(18) в виде:
\begin{equation}
\frac{d^2\psi(z)}{dz^2} - \frac{\psi(z)}{r_D^2} = \frac{4\pi l_B f^{3/2}}{a^3} \sqrt{\psi(z)}.
\end{equation}

Решением этого уравнения является следующий профиль плотности полимера:
\begin{equation}
\rho(z) = \frac{16\pi}{3} \frac{u f^2 r_D^2}{a^5} \sinh^2 \left( \frac{D - z}{4 r_D} \right),
\end{equation}
где толщина адсорбированного слоя $D$ находится из граничного условия~(14):
\begin{equation}
\frac{64\pi}{9} \frac{u f^3 r_D^3}{a^3} \sinh^3\left(\frac{D}{4r_D}\right) \cosh\left(\frac{D}{4r_D}\right) = \sigma a^2.
\end{equation}

Для низких концентраций соли ($r_D \gg D$) профиль плотности полимера имеет параболический вид~\cite{Dobrynin2000a, Dobrynin2000b}:
\begin{equation}
\rho(z) = \frac{\pi}{3} \frac{u f^2 (D - z)^2}{a^5}, \quad \text{($\theta$-растворитель)},
\end{equation}
а толщина слоя $D$ задается выражением:
\begin{equation}
D = \left(\frac{9}{\pi}\right)^{1/3} a^{5/3} u^{-1/3} f^{-1} \sigma^{1/3}, \quad \text{($\theta$-растворитель)}.
\end{equation}

Адсорбированный слой можно представить как составленный из блобов с постепенно возрастающим размером $\xi(z)$~\cite{Dobrynin2000a, Dobrynin2000b}. Число мономеров $g(z)$ в блобе определяется тем, что блобы плотно заполняют пространство, $g(z) \approx \rho(z)\xi^3(z)$, а статистика цепи внутри блоба является гауссовой, $\xi^2(z) \approx a^2 g(z)$. Это дает зависимость размера блоба
\begin{equation}
\xi(z) \approx \frac{a^3}{u f^2 (D - z)^2}, \quad \text{($\theta$-растворитель)}
\end{equation}
и числа мономеров в нем:
\begin{equation}
g(z) \approx \frac{a^4}{u^2 f^4 (D - z)^4}, \quad \text{($\theta$-растворитель)}
\end{equation}
на расстоянии $z$ от поверхности. Такие блобы представляют собой многовалентные участки полиионов с зарядом $q(z) = f g(z)$.

Профиль плотности~(23) не применим непосредственно у заряженной поверхности ($z \approx 0$) и на границе адсорбированного слоя ($z \approx D$) --- в областях с большими градиентами плотности полимера. Вблизи поверхности типичные флуктуации плотности происходят на масштабах порядка корреляционной длины $\xi(z) \approx (\rho(z) a^2)^{-1}$. Таким образом, толщина обедненного слоя у поверхности имеет порядок $\xi(0) \approx a^{1/3} u^{-1/3} \sigma^{-2/3}$. Вкладом этого слоя можно пренебречь, пока корреляционная длина $\xi(0)$ меньше толщины адсорбированного слоя $D$, что выполняется при плотностях поверхностного заряда $\sigma$ больших, чем
\begin{equation}
\sigma_e \approx f / a^2, \quad \text{($\theta$-растворитель)}.
\end{equation}

На внешней границе адсорбирующего слоя трехчастичные взаимодействия $a^6 \rho(z)^3$ становятся порядка конформационной энтропии цепи:
\begin{equation}
\frac{u^3 f^6 (D - z)^6}{a^9} \approx a^6 \rho(z)^3 \approx \frac{a^2}{6} \left( \frac{d\sqrt{\rho(z)}}{dz} \right)^2 \approx \frac{u f^2}{a^3}
\end{equation}
на расстояниях от заряженной поверхности порядка
\begin{equation}
z \approx D - a (u f^2)^{-1/3} \approx D - D_e.
\end{equation}
На этих масштабах необходимо учитывать градиентный член в уравнении~(17). Однако вкладом этого слоя также можно пренебречь, если толщина адсорбированного слоя $D$ превышает размер электростатического блоба $D_e$ (то есть при $\sigma > \sigma_e$).

Интегрирование профиля плотности полимера $\rho(z)$ от $0$ до $D$ дает поверхностный избыток полимера:
\begin{equation}
\Gamma = \int_0^D \rho(z) dz = \frac{3}{8}\frac{\sigma}{f} \left( \frac{2\sinh(y/2) - y}{\sinh^3(y/4)\cosh(y/4)} \right) \approx \frac{\sigma}{f} \left( 1 - \frac{D^2}{20 r_D^2} \right),
\end{equation}
где $y = D / r_D$. Таким образом, при нулевой концентрации соли адсорбированные полиэлектролиты точно компенсируют поверхностный заряд $\sigma$.

Добавление соли уменьшает количество адсорбированного полимера, поскольку ионы соли также участвуют в экранировании поверхностного заряда. Специфическую степенную зависимость поверхностного избытка от концентрации соли можно понять из простых соображений. Типичная избыточная плотность заряда $\delta\rho_{\text{salt}}$ ионов соли в адсорбирующем слое равна
\begin{equation}
\delta\rho_{\text{salt}}(z) \approx \rho_{\text{salt}}^-(z) - \rho_{\text{salt}}^+(z) \approx \rho_{\text{salt}} \psi(z),
\end{equation}
где $\psi(z)$ --- электростатический потенциал на расстоянии $z$ от поверхности. Типичное значение электростатического потенциала $\psi(z)$ в слое можно оценить как $l_B \sigma D$. Умножая избыточную плотность заряда $\delta\rho_{\text{salt}}$ на толщину слоя $D$, получаем поверхностный избыток противионов:
\begin{equation}
\Gamma_{\text{salt}} \approx \delta\rho_{\text{salt}} D \approx l_B \rho_{\text{salt}} \sigma D^2 \approx \sigma \frac{D^2}{r_D^2}.
\end{equation}
Таким образом, часть поверхностного заряда величиной $\sigma D^2 / r_D^2$ экранируется ионами соли, оставляя полиэлектролитным цепям компенсацию лишь доли $\sigma - \sigma D^2 / r_D^2$.

\subsection{Перезарядка поверхности}

Как отмечалось выше, описание адсорбированного слоя в приближении среднего поля перестает работать на расстояниях $z > D - D_e$. На этих масштабах флуктуации плотности полимера $a^{-3}(u f^2)^{1/3}$ становятся больше средней плотности, даваемой формулой~(23). Чтобы учесть перезарядку поверхности адсорбированными полиэлектролитами из-за флуктуаций плотности на границе слоя, разобьем раствор на две области. В области~I толщиной $D$ экранирование обеспечивается полиэлектролитными цепями и ионами соли. В области~II ($z > D$) концентрация полимера равна нулю, а экранирование определяется исключительно ионами соли.

В области~II электростатический потенциал удовлетворяет уравнению
\begin{equation}
\frac{d^2\psi_{\text{II}}(z)}{dz^2} = \frac{\psi_{\text{II}}(z)}{r_D^2}
\end{equation}
с граничным условием
\begin{equation}
\left.\frac{d\psi_{\text{II}}(z)}{dz}\right|_{z=D} = -4\pi l_B \delta\sigma,
\end{equation}
где $\delta\sigma = \sigma - f\Gamma - \Gamma_{\text{salt}}$ --- эффективная поверхностная плотность заряда, величина которой должна быть найдена самосогласованно из сшивки потенциала и его производной на границе $z = D$. Положительное значение $\delta\sigma$ отвечает недокомпенсации поверхностного заряда. Электростатический потенциал в области~II равен:
\begin{equation}
\psi_{\text{II}}(z) = 4\pi l_B r_D \delta\sigma \exp\left( -\frac{z - D}{r_D} \right).
\end{equation}

В области~I потенциал и концентрация полимера связаны системой уравнений:
\begin{equation}
g_e^{-1} = a^6 \rho(z)^2 - f \psi_{\text{I}}(z) - \frac{a^2}{6}\frac{1}{\sqrt{\rho(z)}}\frac{d^2\sqrt{\rho(z)}}{dz^2},
\end{equation}
\begin{equation}
\frac{d^2\psi_{\text{I}}(z)}{dz^2} - \frac{\psi_{\text{I}}(z)}{r_D^2} = 4\pi l_B f \rho(z).
\end{equation}
Оценку эффективной плотности заряда $\delta\sigma$ по порядку величины можно получить, оценив члены в последнем уравнении при $z = D$:
\begin{equation}
g_e^{-1} \approx -4\pi f l_B \delta\sigma r_D - \left. \frac{a^2}{6}\frac{1}{\sqrt{\rho(z)}}\frac{d^2\sqrt{\rho(z)}}{dz^2} \right|_{z=D}.
\end{equation}
Второе слагаемое в правой части~(38) оценивается извлечением квадратного корня из уравнения~(37) и взятием второй производной по $z$. Используя условие непрерывности потенциала, производные $\psi_{\text{I}}(z)$ при $z = D$ выражаются через потенциал $\psi_{\text{II}}(z)$ области~II. После преобразований уравнение~(38) принимает вид:
\begin{equation}
g_e^{-1} \approx -4\pi f l_B \delta\sigma r_D - \frac{a^2}{24 r_D^2}.
\end{equation}
Решая это уравнение относительно $\delta\sigma$, находим:
\begin{equation}
\delta\sigma \approx -\frac{f^{1/3}}{u^{1/3} a r_D} - \frac{a}{f u r_D^3} \approx -\sigma_e \frac{D_e}{r_D} \left( 1 + \frac{D_e^2}{r_D^2} \right), \quad \text{($\theta$-растворитель)}
\end{equation}
для величины перезарядки поверхности адсорбированными цепями. Перекомпенсация поверхностного заряда обусловлена положительным химическим потенциалом мономеров в объеме $g_e^{-1}$. Важно подчеркнуть, что величина перезарядки $\delta\sigma$ универсальна и не зависит от исходной плотности заряда поверхности $\sigma$.

\subsection{Экранирование адсорбирующей поверхности малыми ионами}

Экранирование заряженной поверхности определяется преимущественно полиэлектролитами до тех пор, пока толщина адсорбированного слоя $D$ меньше длины Гуи--Чепмена $\lambda = (2\pi l_B \sigma)^{-1}$, либо плотность поверхностного заряда $\sigma$ меньше некоторого значения $\sigma_{\text{ion}}$:
\begin{equation}
\sigma_{\text{ion}} \approx \frac{f^{3/4}}{a^2 u^{1/2}}, \quad \text{($\theta$-растворитель)}.
\end{equation}
При плотности заряда $\sigma \approx \sigma_{\text{ion}}$ на поверхностный блоб приходится в среднем один заряженный мономер: $f g(0) \approx 1$. При этом концентрация мономеров у поверхности достигает $a^{-3} f^{1/2}$. Дальнейший рост полимерной плотности невыгоден из-за больших энергетических затрат на короткодействующее отталкивание мономер--мономер. При более высоких плотностях заряда ($\sigma > \sigma_{\text{ion}}$) противионы начинают доминировать в экранировании потенциала внутри приповерхностного слоя толщиной $h$. Противион теряет порядка $kT$ трансляционной энтропии при локализации у поверхности в слое толщиной $h$, тогда как увеличение полимерной концентрации выше $a^{-3} f^{1/2}$ потребовало бы энергии отталкивания, существенно превышающей $kT$.

Для получения аналитических выражений разобьем адсорбированный слой на две области. Внутри области~I толщиной $h$ экранирование электрического поля поверхности осуществляется преимущественно противионами. Уравнение Пуассона--Больцмана для этой области имеет вид:
\begin{equation}
\frac{d^2 \delta\psi(z)}{dz^2} = 4\pi l_B \rho_S(h) \exp(\delta\psi(z)),
\end{equation}
где $\rho_S(h)$ --- плотность противионов при $z = h$, а $\delta\psi(z) = \psi(z) - \psi(h)$. Электростатический потенциал в области~I имеет логарифмический вид:
\begin{equation}
\delta\psi(z) = -2\ln\left( 1 + \frac{z - h}{\lambda_2} \right),
\end{equation}
где параметр $\lambda_2 = 1/\sqrt{2\pi l_B \rho_S(h)}$. На заряженной поверхности производная $d\,\delta\psi(z)/dz$ должна быть равна $-4\pi l_B \sigma$. Это требование дает соотношение:
\begin{equation}
\lambda_2 = \lambda + h.
\end{equation}
Распределение противионов в этой области описывается выражением:
\begin{equation}
\rho_S(z) = \frac{1}{2\pi l_B} (\lambda_2 + z - h)^{-2}.
\end{equation}
Профиль плотности полимера в области~I слабо (логарифмически) зависит от расстояния $z$:
\begin{equation}
\rho(z) \approx a^{-3} f^{1/2} \sqrt{\psi(h) - 2\ln\left( 1 + \frac{z - h}{\lambda_2} \right)}.
\end{equation}

Локализованные противионы (внутри области~I) уменьшают поверхностную плотность заряда $\sigma$ до эффективного значения $\sigma_2$, которое задает граничное условие для нелинейного уравнения~(20) в области~II:
\begin{equation}
\left.\frac{d\psi(z)}{dz}\right|_{z=h} = -4\pi l_B \sigma_2.
\end{equation}

В области~II экранирование контролируется адсорбированными полиэлектролитами. Решение уравнения~(20) при низких концентрациях соли по-прежнему дает параболический профиль плотности:
\begin{equation}
\rho(z) \approx \frac{\pi}{3}\frac{u f^2 (D_{\sigma_2} + h - z)^2}{a^5},
\end{equation}
где толщина $D_{\sigma_2}$ области~II определяется формулой~(24) с заменой плотности заряда $\sigma$ на эффективную плотность $\sigma_2$ на границе областей ($z = h$). Граница между двумя областями определяется плоскостью, где локальная концентрация противионов $\rho_S(h)$ сравнивается с концентрацией заряженных мономеров $f\rho(h)$:
\begin{equation}
\frac{1}{2\pi l_B \lambda_2^2} = \frac{\pi}{3}\frac{u f^3 D_{\sigma_2}^2}{a^5}.
\end{equation}
Совместное решение уравнений~(44) и~(49) позволяет найти эффективную поверхностную плотность заряда $\sigma_2$:
\begin{equation}
\sigma_2 = \frac{6^{1/4}}{2\sqrt{\pi}} \frac{f^{3/4}}{u^{1/2} a^2} \approx \sigma_{\text{ion}}.
\end{equation}
Полная толщина адсорбированного полимерного слоя $D$ равна сумме толщины $h$ области~I и толщины $D_{\sigma_2}$ области~II:
\begin{equation}
D = h + D_{\sigma_2} = \frac{4}{6^{1/4}\sqrt{\pi}} \frac{a}{u^{1/2} f^{3/4}} - \lambda.
\end{equation}
При высоких плотностях заряда ($\sigma \gg \sigma_{\text{ion}}$) толщина слоя выходит на насыщение: $D \to \dfrac{4}{6^{1/4}\sqrt{\pi}} \dfrac{a}{u^{1/2} f^{3/4}}$.


\begin{english}
\begin{thebibliography}{99}

\bibitem{Hara1993}
\textit{Polyelectrolytes}, edited by M.~Hara (Marcel Dekker, New York, 1993).

\bibitem{Schmitz1993}
K.~S.~Schmitz, \textit{Macroions in Solution and Colloidal Suspension} (VCH, New York, 1993).

\bibitem{Kawaguchi1992}
M.~Kawaguchi and A.~Takahashi, Adv. Colloid Interface Sci. \textbf{37}, 219 (1992).

\bibitem{Fleer1993}
G.~J.~Fleer, M.~A.~Cohen Stuart, J.~M.~H.~M.~Scheutjens, T.~Cosgrove, and B.~Vincent, \textit{Polymers at Interfaces} (Chapman and Hall, London, 1993).

\bibitem{Bajpai1997}
A.~K.~Bajpai, Prog. Polym. Sci. \textbf{22}, 523 (1997).

\bibitem{VanderSchee1984}
H.~A.~Van der Schee and J.~Lyklema, J. Phys. Chem. \textbf{88}, 6661 (1984).

\bibitem{Evers1986}
O.~A.~Evers, G.~J.~Fleer, J.~M.~H.~M.~Scheutjens, and J.~Lyklema, J. Colloid Interface Sci. \textbf{111}, 446 (1986).

\bibitem{Blaakmeer1990}
J.~Blaakmeer, M.~R.~Böhmer, M.~A.~Cohen Stuart, and G.~J.~Fleer, Macromolecules \textbf{23}, 2301 (1990).

\bibitem{Bohmer1990}
M.~R.~Böhmer, O.~A.~Evers, and J.~M.~H.~M.~Scheutjens, Macromolecules \textbf{23}, 2288 (1990).

\bibitem{Varoqui1991}
R.~Varoqui, A.~Johner, and A.~Elaissari, J. Chem. Phys. \textbf{94}, 6873 (1991).

\bibitem{Varoqui1993}
R.~Varoqui, J. Phys. II \textbf{3}, 1097 (1993).

\bibitem{Borukhov1995}
I.~Borukhov, D.~Andelman, and H.~Orland, Europhys. Lett. \textbf{32}, 499 (1995).

\bibitem{Chatellier1996}
X.~Chatellier and J.-F.~Joanny, J. Phys. II \textbf{6}, 1669 (1996).

\bibitem{deGennes1979}
P.-G. de Gennes, \textit{Scaling Concepts in Polymer Physics} (Cornell University Press, Ithaca, 1979).

\bibitem{Dobrynin1995}
A.~V. Dobrynin, R.~H. Colby, and M.~Rubinstein, Macromolecules \textbf{28}, 1859 (1995).

\bibitem{Varoqui1991}
R.~Varoqui, A.~Johner, and A.~Elaissari, J. Chem. Phys. \textbf{94}, 6873 (1991).

\bibitem{Varoqui1993}
R.~Varoqui, J. Phys. II \textbf{3}, 1097 (1993).

\bibitem{Borukhov1995}
I.~Borukhov, D.~Andelman, and H.~Orland, Europhys. Lett. \textbf{32}, 499 (1995).

\bibitem{Chatellier1996}
X.~Chatellier and J.-F.~Joanny, J. Phys. II \textbf{6}, 1669 (1996).

\bibitem{Borukhov1998}
I.~Borukhov, D.~Andelman, and H.~Orland, Macromolecules \textbf{31}, 1665 (1998).

\bibitem{Joanny1999}
J.-F.~Joanny, Eur. Phys. J. B \textbf{9}, 117 (1999).

\bibitem{deGennes1976}
P.-G. de Gennes, P.~Pincus, R.~M. Velasco, and F.~Brochard, J. Phys. (Paris) \textbf{37}, 1461 (1976).

\bibitem{Israelachvili1985}
J.~N. Israelachvili, \textit{Intermolecular and Surface Forces} (Cornell University Press, Ithaca, 1985).

\bibitem{Dobrynin2000a}
A.~V. Dobrynin, A.~Deshkovskii, and M.~Rubinstein, Phys. Rev. Lett. \textbf{84}, 3101 (2000).

\bibitem{Dobrynin2000b}
A.~V. Dobrynin, A.~Deshkovskii, and M.~Rubinstein, Macromolecules (submitted).


\end{thebibliography}
\end{english}

\end{document}